\section{Introduction}

Recent reasoning-focused language models have renewed interest in test-time scaling, where extra inference-time computation is used to improve answer quality rather than expanding training data or model parameters. The s1 project argues that strong reasoning gains do not necessarily require a complex training stack, large proprietary datasets, or opaque systems. Instead, it combines a compact supervised fine-tuning dataset, called s1K, with a decoding-time intervention named budget forcing \citep{muennighoff2025s1}.

This repository focuses on the budget forcing component because it is narrow, inspectable, and experimentally accessible. In contrast to larger reinforcement-learning-based systems such as DeepSeek-R1 \citep{deepseek2025r1} or large-scale inference-oriented models such as Kimi k1.5 \citep{kimi2025k15}, budget forcing modifies generation at inference time by either extending or truncating the model's reasoning trace. That makes it a suitable target for a graduate reproduction project with limited compute.

The present work has two goals. First, it audits and executes the existing implementation in this repository to establish whether the reported budget-forcing behavior can be reproduced under a modest experimental setup. Second, it analyzes the mechanism and its limits through the lens of reproducibility: how much of the observed behavior depends on model choice, trigger design, benchmark selection, and hardware availability.

The main contributions of this report are as follows. We provide a structured summary of the s1 method and its evaluation metrics, instantiate an end-to-end experiment pipeline for budget forcing on math reasoning tasks, and document a paper-style reproduction workflow that separates stable narrative content from result-dependent content. This separation is useful because the repository was already close to feature-complete, while the missing work was concentrated in execution, reporting, and integration.

The remainder of the paper is organized as follows. Section~\ref{sec:method} summarizes the s1 method and the local implementation. Section~\ref{sec:metrics} explains the three evaluation metrics used to analyze test-time scaling. Section~\ref{sec:analysis} discusses strengths, limitations, and research gaps. Section~\ref{sec:related} connects budget forcing to adjacent work on reasoning models and inference-time augmentation. Section~\ref{sec:experiments} documents the experimental setup and provides a template for presenting the observed results. Section~\ref{sec:conclusion} concludes with practical lessons and future directions.

Test-time scaling has become a central topic in recent reasoning-focused large language model research because it offers a path to stronger performance without changing the base pretraining corpus at inference time. Instead of relying only on larger models or more supervised data, test-time scaling methods spend additional compute during generation, typically through longer reasoning traces, sampling, verification, or structured control policies. Recent public discussion around reasoning models has intensified because systems such as OpenAI's o1 demonstrated strong reasoning behavior without fully disclosing the recipe behind that behavior.

The s1 framework is important in this context because it argues that competitive reasoning behavior can emerge from a comparatively simple recipe: a small, carefully curated dataset of reasoning traces and a decoding-time intervention called budget forcing \citep{muennighoff2025s1}. Rather than treating reasoning depth as a fixed model property, budget forcing manipulates the generation process directly. When the model tries to stop too early, the decoder can suppress termination and append a trigger phrase such as ``Wait'', effectively encouraging more internal computation. Conversely, when a hard budget is reached, the decoder can force a transition to the final answer segment.

This project examines that claim from a reproduction perspective. Our repository focuses on the decoding-time mechanism rather than a full-scale retraining pipeline. The implemented code already contains a budget forcing decoder, an evaluation entrypoint, basic scripts for baseline and budget forcing sweeps, and notebooks for reproduction and trigger ablation. The main research task now is to turn these components into validated experimental artifacts and a coherent analysis.

The report makes three practical contributions. First, it explains the s1 method and the role of budget forcing in accessible terms. Second, it documents the current repository implementation and its evaluation design. Third, it frames reproduction as an engineering problem with explicit environmental constraints, especially for local execution on non-CUDA hardware. This is important because negative or partial reproduction results can still be scientifically useful when they clarify what parts of a claimed method are robust and what parts depend on hidden infrastructure assumptions.

The remainder of this report is organized as follows. Section~\ref{sec:method} summarizes the s1 pipeline and the implemented budget forcing mechanism. Section~\ref{sec:metrics} explains the control, scaling, and performance metrics used to judge test-time scaling behavior. Section~\ref{sec:analysis} discusses the strengths, limits, and research gaps that motivate our reproduction. Section~\ref{sec:related} positions the work relative to adjacent reasoning and inference-time augmentation methods. Section~\ref{sec:experiments} describes the repository experiments and records where result integration depends on successful execution. Section~\ref{sec:conclusion} concludes with key observations and next steps.